\section{Method}\label{sec:method}

The hockey environment's fully continuous action space makes discrete-action methods such as
\ac{DQN}~\cite{mnih2015human:DQN} inapplicable, and on-policy algorithms such as
\ac{PPO}~\cite{SchulmanEtAl2017:PPO} less attractive due to poor sample efficiency.
We therefore adopt two off-policy, continuous-control algorithms that share a common backbone
yet represent complementary design philosophies: \ac{TD3}~\cite{fujimoto2018:TD3} optimizes a
deterministic policy with explicit bias-correction mechanisms, while \ac{SAC}~\cite{haarnoja2018:SOFT}
optimizes a stochastic policy through entropy maximization.
Both maintain a replay buffer $\mathcal{D}$ of transitions $(s_t, a_t, r_t, s_{t+1}, d_t)$,
train an actor $\pi$ and one or more critics $Q_\theta$ on minibatches from $\mathcal{D}$, and
stabilize learning via slowly updated target networks through Polyak averaging,
\begin{equation}
  \theta' \leftarrow \tau\,\theta + (1-\tau)\,\theta', \qquad
  \phi'   \leftarrow \tau\,\phi  + (1-\tau)\,\phi',
  \label{eq:polyak}
\end{equation}
where $\tau \ll 1$ ensures the target parameters $\theta'$ and $\phi'$ track their online
counterparts $\theta$ and $\phi$ slowly, providing a stable and consistent learning target
throughout training.
The individual algorithms are described in Sections~\ref{sec:td3} and~\ref{sec:sac},
further improvements to both are presented in Section~\ref{sec:improvements}.


\subsection{Twin Delayed Deep Deterministic Policy Gradient}
\label{sec:td3}

\ac{TD3}~\cite{fujimoto2018:TD3} extends \ac{DDPG}~\cite{lillicrap2015:DDPG} by adressing two fundamental failure modes of determenistic actor-critic methods:
overestimation bias in the critic and instability in polciy updates.
The actor encodes a \emph{deterministic} policy $a = \pi_\phi(s)$ and is trained to maximize expected discounted return via the determenistic policy gradient theorem~\cite{silver2014:DPG}:
\begin{equation}
  \nabla_\phi J(\phi)
    \approx \mathbb{E}_{s \sim \mathcal{D}}\!\left[
      \nabla_a Q_{\theta}(s,a)\big|_{a=\pi_\phi(s)}\,
      \nabla_\phi \pi_\phi(s)
    \right].
\end{equation}
In actor-critic methods, the policy is updated greedily with respect to the critic,
meaning any overestimation in the value function is directly exploited by the actor,
leading to suboptimal or diverging behaviour~\cite{fujimoto2018:TD3}.

\ac{TD3} adresses this with three interacting mechanisms.

\vspace{4pt}

\textbf{Clipped double Q-learning}

Two independent critics $Q_{\theta_1}$ and $Q_{\theta_2}$ are trained in parallel but share no parameters.
When computing Bellman targets, the minimum of the two target critics is used:
\begin{equation}
  y = r + \gamma(1-d)\min_{i=1,2} Q_{\theta_i'}(s', \tilde{a})
  \label{eq:td3target}
\end{equation}
where $d$ indicates episode termination and $\theta_i$ are the target network parameters updated via~\eqref{eq:polyak}.
Taking the minimum introduces a conservative bias that counteracts the systematic overastimation inherent in maximizing over noisy function approximators.
Both critics are updated independently by minimizing the mean squarred Bellman error $(y - Q_{\theta_i}(s, a))^2$ over minibatches from $D$.

Since gradients flow independently through each critic, the two networks diverge slightly,
preventing the correlated overestimation that arises when a single critic is used.

\vspace{4pt}

\textbf{Target policy smoothing}

To prevent the actor from exploiting narrow, misleading peaks in the value function, 
the target action is perturbed with clipped Gaussian noise before evaluating the Bellman target:
\begin{equation}
  \tilde{a} = \operatorname{clip}\!\left(\pi_{\phi'}(s') + \epsilon,\;a_{\min},\;a_{\max}\right),
  \quad \epsilon \sim \operatorname{clip}\!\left(\mathcal{N}(0,\sigma),\;-c,\;c\right).
\end{equation}
This encourages the critics to learn a smoother value landscape, making individual erroneous peaks less exploitable.
The noise standard deviation $\sigma$ and clip range $c$ are hyperparameters that control the effective radius of this smoothing.

\vspace{4pt}

\textbf{Delayed policy updates}

Even with conservative targets and smoothed backups, updating the actor at every critic step introduces high variance into the policy gradient:
the critic has not yet fully integrated the latest target information before its gradient is used to move the policy.
\ac{TD3} therefore updates the actor every $K$ critic steps, typically $K=2$.
This delay ensures the critics are in a more stable and accurate state when the actor gradient is computed,
substantially reducing variance of policy updates and preventing the actor from over-adapting to transient critic errors~\cite{fujimoto2018:TD3}.
Target networks for both actor and critics are updated via the shared Polyak rule~\eqref{eq:polyak}, 
ensuring the learning targets change slowly and the training signal remains consistent across consecutive gradient steps.

Togehter, these three mechanisms directy counteract the core failure modes of deteremenistic actor-critic methods:
clipped double Q-learning reduces overestimation bias, target policy smoothing regularizes the value function against sharp local peaks,
and delayed updates reduce gradient variance during policy optimization.

\subsection{Soft Actor-Critic}
\label{sec:sac}

\ac{SAC} is an off-policy actor-critic algorithm designed for optimizing \emph{stochastic} policies in
continuous action spaces.
Our implementation follows the framework introduced by~\cite{haarnoja2018:SOFT}, which augments the
standard maximum reward objective with an entropy maximization term~\cite{ziebart2010}.
This approach encourages wider exploration and improves robustness by seeking a policy that
maximizes the trade-off between expected return and entropy:
\begin{equation}
  J(\pi) = \sum_{t=0}^{T}
    \mathbb{E}_{(\mathbf{s}_t,\mathbf{a}_t)\sim\rho_\pi}\!\left[
      r(\mathbf{s}_t,\mathbf{a}_t) + \alpha\,\mathcal{H}\!\left(\pi(\cdot|\mathbf{s}_t)\right)
    \right]
\end{equation}
where $\rho_\pi$ denotes the state-action distribution induced by policy $\pi$ and $\mathcal{H}$
represents the entropy of the policy at state $\mathbf{s}_t$.

The \ac{SAC} algorithm consists of an actor, critic, and an entropy regularization term governed
by a temperature parameter $\alpha$.
The actor is parameterized by a neural network that defines a Gaussian distribution over the
continuous action space by estimating the mean $\mu_\theta(s)$ and standard deviation
$\sigma_\theta(s)$ in its output layer.

To ensure differentiability of the sampling process with respect to the policy parameters, the
reparameterization trick is applied~\cite{haarnoja2018:SOFT, kingma2014auto}.
Instead of sampling actions directly from the Gaussian distribution, actions are expressed as
$a = \mu_\theta(s) + \sigma_\theta(s) \odot \epsilon$ where $\epsilon \sim \mathcal{N}(0, I)$ is
a standard normal random variable.

The sampled actions are passed through a $\tanh$ function and scaled to match the action bounds
of the environment.
To correctly account for this squashing transformation in the entropy term, the log-probability
is adjusted using the determinant of the Jacobian of the $\tanh$ transformation.
The actor is trained on the maximum entropy objective~\cite{haarnoja2018:SOFT}:
\begin{equation}
  J(\pi_\phi) =
    \mathbb{E}_{(\mathbf{s}_t \sim \mathcal{D})}\!\left[
      \mathbb{E}_{(\mathbf{a}_t) \sim \pi_\phi}\!\left[
        \alpha \log \pi_\phi(a_t|s_t) - Q_\theta(s_t, a_t)
      \right]
    \right]
\end{equation}
where for each state sampled from the replay buffer $\mathcal{D}$, an action is sampled from the
actor and the critic estimates the value for each state-action pair.

As presented in the follow-up work~\cite{haarnoja2018algorithms}, we employ two Q-networks for
the critic instead of a single Q-network and a value network as originally proposed.
The soft value function can be implicitly derived from the soft Q-function by combining the
Q-values with the policy entropy term.

During training, two independent Q-networks estimate the value of a state-action pair.
To mitigate overestimation bias, soft policy evaluation is performed using the minimum of the two
Q-values, which provides a more conservative value estimate~\cite{fujimoto2018:TD3}.
The critic is trained using a mean squared Bellman error loss:
\begin{equation}
  J_Q(\theta_i) = \mathbb{E}_{\mathcal{D}}\!\left[
    \left(Q_{\theta_i}(s_t,a_t) - y_t\right)^2
  \right]
\end{equation}
with target value
\begin{equation}
  y_t = r_t + \gamma\!\left(
    \min_{j=1,2} Q_{\bar{\theta}_j}(s_{t+1},a_{t+1})
    - \alpha \log \pi_\phi(a_{t+1}|s_{t+1})
  \right).
\end{equation}
The target networks $\bar{\theta}_j$ are updated using soft updates with factor $\tau$, which
improves training stability by preventing rapid divergence between online and target critics.

Furthermore, automatic parameter tuning is implemented as introduced in the follow-up
paper~\cite{haarnoja2018algorithms}.
The temperature parameter $\alpha$ is optimized to match a target policy entropy, allowing the
agent to automatically balance exploration and exploitation.
The temperature is updated by minimizing
\begin{equation}
  J(\alpha) = \mathbb{E}_{a_t \sim \pi_\phi}\!\left[
    -\alpha\left(\log \pi_\phi(a_t|s_t) + \bar{\mathcal{H}}\right)
  \right],
\end{equation}
where $\bar{\mathcal{H}}$ denotes the target entropy, typically set to
$-\operatorname{dim}(\mathcal{A})$ as recommended in~\cite{haarnoja2018algorithms}.
The parameter $\alpha$ is optimized using a gradient-based method.




\subsection{Improvements}
\label{sec:improvements}

We implemented a framework based on YAML configuration files to manage and persist training runs,
enabling reproducible experiments and systematic tracking of hyperparameter changes.
All modifications described below are applied on top of both the \ac{TD3} and \ac{SAC} base implementations.


\subsubsection*{Replay Buffer}
Since both implemented algorithms are off-policy, training relies on data stored in a replay
buffer. We implemented standard \ac{ER}, which improves data efficiency by reusing past 
interactions and reducing correlations between subsequent samples. Transitions are stored as 
tuples $(s_t, a_t, r_t, s_{t+1})$ in a finite buffer; during training, batches are sampled 
uniformly from the buffer. However, uniform sampling treats all transitions equally, while 
transitions with large \ac{TD} errors may provide stronger learning signal. To address this 
limitation, we additionally implement \ac{PER}.

\ac{PER} assigns higher sampling probabilities to transitions with larger \ac{TD} 
errors~\cite{Schaul2015PrioritizedER}. Each transition $i$ is associated with a priority 
value $p_i = |\delta_i| + \varepsilon$, where $\delta_i$ denotes the \ac{TD} error and 
$\varepsilon > 0$ ensures non-zero probability. Transitions are sampled according to
\begin{equation}   
  P(i) = \frac{p_i^\alpha}{\sum_k p_k^\alpha}, 
\end{equation}
where $\alpha$ determines the degree of prioritization ($\alpha = 0$ recovers uniform 
sampling). Prioritized sampling introduces a bias into the gradient estimates, corrected 
via importance-sampling weights:
\begin{equation}   
  w_i = \left(\frac{1}{N} \cdot \frac{1}{P(i)}\right)^\beta, 
\end{equation}
where $N$ is the buffer size and $\beta$ controls the strength of bias 
correction~\cite{Schaul2015PrioritizedER}. To ensure asymptotically unbiased updates, 
$\beta$ is annealed linearly to $1$ over a configurable fraction of total training steps. 
Efficient sampling and priority updates are implemented using a binary SumTree data 
structure, enabling both operations in $\mathcal{O}(\log N)$ time.

\subsubsection*{Pink Noise}
Exploration in off-policy algorithms is commonly achieved by injecting noise into the 
actor's actions during training. The standard choice is white noise $\mathcal{N}(0, \sigma)$, 
which is temporally uncorrelated, meaning each noise sample is drawn independently at 
every timestep, leading to jittery exploration with no coherence across consecutive actions. 
Eberhard et al. introduced colored noise as a principled alternative, parameterized by a 
spectral exponent $\beta$ that controls the degree of temporal correlation~\cite{eberhard2023pink}. 
Increasing $\beta$ yields smoother trajectories where the noise is correlated with past samples, 
encouraging the agent to commit to exploratory directions over multiple consecutive steps. 
Pink noise ($\beta = 1$) was shown to strike the best balance between exploitation and 
exploration across a range of continuous control tasks~\cite{eberhard2023pink}.

Colored noise is generated by decomposing a white noise signal into its frequency components 
via the Fast Fourier Transform, scaling each component $f$ by $f^{-\beta/2}$, and transforming 
the result back into the time domain, where frequency refers to how rapidly the noise changes 
over consecutive timesteps. Suppressing high frequencies and amplifying low frequencies 
produces a noise sequence that varies slowly over time, resulting in temporally correlated 
action perturbations across the agent's trajectory. We apply pink noise as the exploration 
noise for \ac{TD3}, replacing the standard Gaussian perturbation. Since \ac{SAC} learns a 
stochastic policy whose entropy is explicitly regularized through $\alpha$, it does not rely 
on external exploration noise and pink noise is therefore not applicable.

\subsubsection*{Reward Shaping}
To enhance learning in the hockey environment, we extend the default reward formulation 
with additional shaping terms. The environment provides a sparse reward for winning, a 
penalty for losing, and optional signals for puck contact and puck velocity direction, 
which are not enabled by default. The total reward is given by $r = \sum_{i} w_i \cdot r_i$,
where each $r_i$ denotes a reward component and $w_i \in \mathbb{R}$ its corresponding weight.

In early training stages, the agent often required a significant number of timesteps 
before initiating contact with the puck. To accelerate early learning, we introduce a 
time-penalty term $r_{\text{passive}}$, applied when the puck remains stationary on the 
agent's side of the field beyond the first $10\%$ of episode steps,
\begin{equation}
    r_{\text{passive}}(t) = \begin{cases} -c & \text{if } x_{\text{puck}} \leq 0, \; 
    \|v_{\text{puck}}\| < 0.1, \text{ and } t > 0.1 \cdot T_{\max}, \\ 
    0 & \text{otherwise,} \end{cases}
\end{equation}
where $c > 0$ is a configurable penalty coefficient. To encourage defensive positioning, 
we additionally reward the agent for staying close to its own goal when the puck is on 
the opponent's side,
\begin{equation}
    r_{\text{defense}} = \begin{cases} \dfrac{1}{1 + \|p_{\text{agent}} - g\|_2} 
    & \text{if } x_{\text{puck}} > 0, \\ 0 & \text{otherwise,}\end{cases}
\end{equation}
where $g$ denotes the agent's goal center. All shaping components are individually 
configurable via the YAML configuration file and can be freely weighted or disabled.


\subsubsection*{Self-Play}


\begin{figure}[htbp]
\centering
\includegraphics[width=0.85\textwidth]{./figures/selfplay_diagram.pdf}
\caption{Self-play opponent pool workflow (AlphaZero-inspired). Future extension for TD3.}
\label{fig:selfplay}
\end{figure}
